\begin{figure}[ht]
	\centering
	\begin{tcolorbox}[sharp corners,colback=white,boxrule=0.3mm,left=1mm]
		\begin{subfigure}{\textwidth}
			\centering
			\begin{tabular}{|c|c|}
				\hline
				$\otpct_1^\rev =  \varepsilon \cdot \pker.\enc(\revpk^1_1, m) $
				 & $ct^\mat_1 =  \pkem.\enc(\matpk^{2}_1, \varepsilon \cdot \alpha_1) $ \\
				\hline
				$\otpct_2^\rev= \alpha_1 \cdot \pker.\enc(\revpk^4_2, m ) $
				 & $ct^\mat_2 =  \pkem.\enc(\matpk^5_2, \alpha_2 \cdot \alpha_1) $      \\
				\hline
				$\otpct_3^\rev=  \alpha_2 \cdot \pker.\enc(\revpk^2_3, m) $
				 & $ct^\mat_3 =  \pkem.\enc(\matpk^3_3, \alpha_3 \cdot \alpha_2) $      \\
				\hline
				$\otpct_4^\rev =  \alpha_3 \cdot \pker.\enc(\revpk^8_4, m ) $
				 & $\bot$                                                               \\
				\hline
			\end{tabular}
			\caption{Example $\escrow$ for $\uval = 1427$.}
		\end{subfigure}

		\vspace{1.5em}

		\begin{subfigure}{\textwidth}
			\centering
			\begin{tabular}{|c|c|}
				\hline
				$\textcolor{red}{\bot =  \pker.\dec(\revsk^j_1, \epsilon^{-1} \cdot \otpct^\rev_1)} $
				 & $\textcolor{blue}{\tilde{\otpk}_1 =  \epsilon^{-1} \cdot \pkem.\dec(\matsk^{2}_1, ct^\mat_1) =\alpha_1} $             \\
				\hline
				$\textcolor{red}{\bot =  \pker.\dec(\revsk^j_{2}, (\tilde{\otpk}_1)^{-1} \cdot \otpct^{\rev}_2)}$
				 & $ \textcolor{blue}{\tilde{\otpk}_2 =  (\tilde{\otpk}_1)^{-1} \cdot \pkem.\dec(\matsk^5_2, ct^\mat_2)   =\alpha_2}$    \\
				\hline
				$\textcolor{red}{\bot =  \pker.\dec(\revsk^j_3, (\tilde{\otpk}_2)^{-1} \cdot \otpct^{\rev}_3)} $
				 & $\textcolor{red}{\tilde{\otpk}_3 =  (\tilde{\otpk}_2)^{-1} \cdot \pkem.\dec(\matsk^9_3, ct^\mat_3)  \neq \alpha_3 } $ \\
				\hline
				$\textcolor{red}{\bot =  \pker.\dec(\revsk^j_4, (\tilde{\otpk}_3)^{-1} \cdot \otpct^{\rev}_4)} $
				 &                                                                                                                       \\
				\hline
			\end{tabular}
			\caption{Example failed $\dec$ for $\uval = 1427 < \thr = 1486$. \textcolor{red}{Red} indicates a failed decryption, while \textcolor{blue}{Blue} indicates a successful one. Note that the decryptor cannot tell whether a decrypted $\tilde{\otpk}_i$ is correct.}
		\end{subfigure}
	\end{tcolorbox}
	\caption{An example of an $\escrow$ and failed $\dec$ for a value $\uval < \thr$.}
	\label{fig:failed_dec_example}
\end{figure}

